\documentclass[11pt]{article}

\usepackage[margin=1in]{geometry}
\usepackage[utf8]{inputenc}
\usepackage[T1]{fontenc}
\usepackage{lmodern}
\usepackage{newunicodechar}
\usepackage[round,authoryear]{natbib}
\usepackage{booktabs}
\usepackage{tabularx}
\newcolumntype{L}{>{\raggedright\arraybackslash}X}
\newcolumntype{R}{>{\raggedleft\arraybackslash}X}
\setlength{\emergencystretch}{3em}
\usepackage{array}
\usepackage{amsmath}
\usepackage{amssymb}
\usepackage{graphicx}
\usepackage[hidelinks]{hyperref}

% Unicode characters used throughout the source text, mapped to their
% standard LaTeX renderings so the running prose can be transcribed
% without manually converting every dash, arrow, or symbol.
\newunicodechar{—}{\textemdash}
\newunicodechar{–}{\textendash}
\newunicodechar{−}{\ensuremath{-}}
\newunicodechar{→}{\ensuremath{\rightarrow}}
\newunicodechar{Δ}{\ensuremath{\Delta}}
\newunicodechar{≥}{\ensuremath{\geq}}
\newunicodechar{≤}{\ensuremath{\leq}}
\newunicodechar{×}{\ensuremath{\times}}
\newunicodechar{·}{\ensuremath{\cdot}}
\newunicodechar{∈}{\ensuremath{\in}}
\newunicodechar{²}{\ensuremath{^2}}
\newunicodechar{ρ}{\ensuremath{\rho}}

\title{Trigger Specificity, Not Correction Precision: A Negative-Results Study of No-Oracle Drift Correction for Frozen SAM2 on SA-V}
\author{}
\date{}

\begin{document}

\maketitle

\begin{abstract}
We study whether a frozen SAM2.1 (Hiera-Tiny) video segmentation model can be improved at inference time, without fine-tuning and without access to ground truth, by detecting and correcting propagation drift — the loss or misassignment of object identity during memory-based mask propagation. On the official SA-V benchmark, a validated baseline prompting policy (\texttt{box\_centroid}) scores J\&F = 72.1 on val and 74.4 on test. A motion-prior drift-correction mechanism (``P5'') that flags suspected drift from area-ratio and position-jump signals and corrects it with a fresh box re-prompt centered on a constant-velocity extrapolated position achieves striking results on a ground-truth-preselected subset of 30 objects known to genuinely drift: +15.5 percentage points J\&F over the objscore no-oracle baseline. Applied without selection to the full, unfiltered 293-object val population, however, the same mechanism is net-harmful: J\&F falls to 68.7, 3.4 points below the plain baseline, with 178 objects regressing against only 44 improving. Twelve further mitigation attempts spanning frame-level gating, object-level circuit breakers, feature-based pre-selection, and correction softening fail to close this gap; the best variant reaches only 0.8 points below baseline on val and 1.8 points below on test. We report these results as a diagnostic negative-results study: correction precision on true positives is not the binding constraint on deployability — trigger specificity under a low positive-class base rate is. We situate the finding relative to imbalanced-classification theory and the test-time-adaptation literature, and we outline untested directions the results motivate without claiming to have tried them.
\end{abstract}

\section{Introduction}

Promptable video object segmentation models built on the Segment Anything paradigm \citep{kirillov2023,ravi2024} let a user specify a target once, typically on the first frame, and propagate a mask through the rest of a video using a learned memory mechanism. This design is attractive precisely because it needs so little interaction: \citet{ravi2024} report that SAM2 requires roughly three times fewer user interactions than prior video segmentation approaches. That efficiency claim rests on an implicit assumption — that a single prompt, propagated by a frozen model's memory, stays anchored to the correct object for the length of the video. When it does not, the practical remedy is either to fine-tune the model on harder data, which is costly and checkpoint-specific, or to find some inference-time way to detect and repair drift without touching the weights and without access to ground truth. The second option is the one this paper studies.

We work with SAM2.1 Hiera-Tiny, frozen throughout, evaluated on the official SA-V benchmark \citep{metaai2024} using the official \texttt{J}, \texttt{F}, \texttt{J\&F} evaluator. A validated baseline prompting policy, \texttt{box\_centroid} — a tight bounding box around the ground-truth mask on the first annotated frame plus a centroid point — scores J\&F = 72.1 on SA-V val (293 objects, 155 videos) and 74.4 on SA-V test (278 objects, 150 videos), after transferring from a J\&F-Mean of 0.8235 on full DAVIS 2017 val. Failure-mode analysis of the val split shows that the dominant cause of low-scoring objects is not a bad initial prompt: 83\% of objects with J\&F below 40 had a correct start (prompt-frame IoU $\geq$ 0.3) but degraded during propagation. An oracle experiment — re-prompting every fifth annotated frame from the ground-truth mask — recovers this population from J\&F = 26.1 to 84.3, a 58.2-point gain, confirming that the degradation is a recoverable memory-persistence problem rather than a perceptual limit of the model.

The question this paper answers is whether that 58.2-point ceiling is reachable without ground truth. We designed a motion-prior drift-correction mechanism, P5, that detects drift from two self-consistency signals (a mask-area ratio against a running baseline, and a position jump against a constant-velocity extrapolation of the object's own trajectory) and corrects it by re-prompting with a fresh bounding box centered on the extrapolated position, rather than by reinjecting the model's own possibly-corrupted mask. On a ground-truth-preselected subset of 30 objects independently identified as genuine propagation-drift failures, fully calibrated P5 reaches J\&F = 42.3, a 15.5-point gain over the no-oracle baseline it replaces in the deployed comparison (27.4, the strongest of the prior no-oracle attempts, is 8.5 points below P5's untuned result) and the best result of the core calibration sweep (a later refinement described in Section 4.5 raises this slightly further, to 42.6).

The paper's central finding is that this result does not survive contact with an unselected population. When the same mechanism is applied to every object in the val split, with no prior knowledge of which objects actually need correction, J\&F drops to 68.7 — 3.4 points \emph{below} the plain uncorrected baseline, with 178 objects regressing against 44 improving. Twelve further mitigation attempts, organized into four families (frame-level gating, object-level circuit breakers, feature-based pre-selection, and correction softening), narrow but never close this gap; the best configuration, a stack of an object-level correction-rate circuit breaker with a frame-level AND-gate, still finishes 0.8 points below baseline on val and 1.8 points below on test.

We present this explicitly as a negative-results and diagnostic paper, not as a paper reporting an improved system. Its contribution is to show, with a full quantitative accounting rather than a curated demonstration, that a correction mechanism's precision on the cases it was designed for is close to irrelevant to whether it is safe to deploy: what determines deployability is the mechanism's \emph{specificity} — how rarely it fires on objects that did not need correction in the first place — under a positive-class base rate of roughly 10\%. Section 2 places this framing against prior work on video object segmentation memory, tracking drift, test-time self-correction, and imbalanced detection. Section 3 describes the model, benchmark, and every mechanism tried. Section 4 reports the full quantitative results, with the curated-subset-versus-full-population contrast as the headline comparison. Section 5 discusses why object-level mechanisms outperformed frame-level ones and why pre-selection and correction-softening both failed. Sections 6–8 cover limitations, literature-derived (untested) future directions, and conclusions.

This paper's contributions are: (1) a methodological point, illustrated independently by both P2 and P5, that a correction mechanism's measured effect on a ground-truth-preselected curated subset can substantially overstate — or even reverse the sign of — its effect on the full, unselected population it would need to run on in deployment; (2) a mechanistic taxonomy explaining why four structurally different mitigation families (frame-level gating, object-level circuit breakers, feature-based pre-selection, and correction softening) each fail for a distinct, diagnosable reason rather than one shared generic cause (Section 5); and (3) release of all code implementing the baseline policies, P5, and all twelve mitigation variants, together with the per-object result CSVs underlying every table in Section 4 (Data Availability Statement).

\section{Related Work}

\subsection{SAM/SAM2 and promptable video object segmentation}

This project operates entirely inside the paradigm SAM established \citep{kirillov2023} and SAM2 extended to video \citep{ravi2024}: prompt-to-mask on a frozen model, with no fine-tuning. The \texttt{box\_centroid} and largest-connected-component prompting policies used here are direct applications of SAM's point/box interface, and the val/test splits (293 and 278 objects respectively) are the official SA-V partitions released alongside SAM2 \citep{metaai2024}. SAM2's own stated efficiency claim — roughly three times fewer interactions than prior interactive VOS methods — describes exactly the regime this paper stress-tests: not whether SAM2 can be prompted well on one frame, but whether a frozen SAM2's own propagation can be corrected, without oracle access, once it drifts. The finding that 83\% of SA-V val failures are propagation-drift with a correct start is a project-original decomposition, but it converges with an independent diagnosis of SAM2's own memory-buffer behavior: \citet{ding2024}, in SAM2Long, describe an ``error accumulation'' pattern in which an erroneous mask written to memory propagates forward and, once SAM2's fixed memory buffer eventually loses the original anchor frame, allows drift onto a visually similar distractor. That two independent analyses — this project's oracle-versus-no-oracle ceiling experiment and SAM2Long's published diagnosis — arrive at the same mechanism from different angles is a genuine confirmation rather than shared vocabulary.

\subsection{Tracking drift, re-identification, and memory-based VOS}

The memory-based VOS lineage SAM2 descends from illustrates a progression of architectural fixes to exactly the problem studied here. STM \citep{oh2019} established dense feature-space memory matching. XMem \citep{cheng2022xmem} added a three-tier sensory/working/long-term memory specifically so long videos do not lose their anchor frame. Cutie \citep{cheng2023cutie} moved from pixel-level to object-level memory reading because pixel-level matching is fragile to visually similar distractors. This project's capstone finding — that a correction mechanism triggered on a frozen model's own propagation signal is net-harmful when applied unselectively — is a symptom of exactly the distractor-fragility problem Cutie's architecture targets, but the frozen-model, inference-only design used here cannot draw on an architectural fix.

Classical visual tracking offers an independent lineage on the same problem. DaSiamRPN \citep{zhu2018} is the earliest precedent for the specific failure P5 targets: ``semantic distractors,'' visually near-identical objects that an appearance-only signal cannot discriminate. Cases in this project's own data — \texttt{sav\_022396} (roughly 15 near-identical balloons) and \texttt{sav\_035221} (two visually identical shoe soles) — are direct instances of this failure mode. Where P5 differs from DaSiamRPN is that its distractor-awareness is purely positional (constant-velocity centroid extrapolation) rather than learned and appearance-based, a choice made specifically because appearance is uninformative between visually identical instances. Siam R-CNN \citep{voigtlaender2020} adds a further relevant precedent: explicit re-detection and tracking of distractor objects alongside the target, resolved through dynamic programming over multiple hypotheses through time, structurally different from — and potentially complementary to — this project's single-hypothesis, per-trigger reprompt.

DEVA \citep{cheng2023deva} offers a related architectural contrast to this project's own single-reference, single-direction attempts to stabilize correction: it decouples task-specific segmentation from a reusable temporal-propagation module and builds a coherent track through bi-directional propagation, fusing segmentation hypotheses from multiple frames rather than committing to one forward-only chain. This is structurally close to what this project's own D3 (a sliding correction-rate window) and consensus-of-K reference-mask variant gesture toward, but DEVA achieves it through a trained bi-directional fusion model rather than a heuristic layered on a frozen video predictor — the structural reason it can resolve ambiguity that this project's single-direction, single-reference heuristics could not.

Two SAM2-specific follow-ups postdate and directly address the mechanism diagnosed independently here. DAM4SAM \citep{videnovic2024} adds a distractor-aware memory with an introspection-based update — deciding whether to trust and write a frame's result \emph{before} committing it to memory — architecturally distinct from, but close in spirit to, this project's B1 (a post-hoc plausibility check on a correction's output) and A2 (an object-level rate circuit breaker). This project later tested DAM4SAM's underlying introspective signal in isolation — the divergence between SAM2's own primary and best-alternative candidate mask, which the tracker computes at every step but normally discards — as a static per-object pre-selection feature (Section 4.5) and found it carries no more predictive signal than E1's trajectory-based features (R² = 0.028 versus E1's 0.026). DAM4SAM's actual mechanism is a different kind of intervention: it acts continuously, frame-by-frame, on what enters a persistent memory bank, rather than as a one-time decision of whether to run correction on an object at all. A later, separate follow-up outside this paper's twelve mitigation attempts subsequently tested that dynamic, memory-level role directly — skipping the write of a suspected-distractor frame into SAM2's non-conditioning memory bank, applied to plain, uncorrected propagation with no reprompt correction at all — and it did not survive contact with the full validation population: a curated n=30 known-drift pilot cleared its pre-committed gate by a wide margin (threshold 0.2, J\&F=39.4 against a 26.1 baseline), but the same threshold produced a net regression on the full 293-object population (J\&F=71.3 against a 72.1 baseline, 0.8 points below), the same curated-subset-versus-full-population pattern this paper documents for P5 (Section 5). The more ambitious half of DAM4SAM's mechanism — injecting a persistent distractor-resolving anchor frame into SAM2's memory-attention window outside its normal recency schedule — remains only partly tested: the original concern that SAM2's learned positional-encoding table has no entry for an arbitrary out-of-schedule distance has since been addressed at the design stage (a proposed route via SAM2's separate, continuously-encoded object-pointer channel, a code path the frozen model already exercises on out-of-range distances for its own conditioning-frame pointer on long videos); a narrower variant has since been implemented and piloted — it swaps the content of an anchor object pointer from the last confirmed-clean frame into the oldest slot of SAM2's object-pointer window, but keeps that slot's fixed positional distance — and it did not pass its pre-committed pilot gate (J\&F=27.8 vs. baseline 26.1 on the n=30 known-drift subset, gate 29.1), while injection at the anchor frame's true, potentially much larger distance remains untested (Section 7). SAM2Long \citep{ding2024} maintains multiple candidate propagation pathways continuously and prunes by cumulative score, rather than committing to a binary detect-then-hard-reprompt decision as P5 does. SAM2Long was not implemented or tested in this project and is flagged as an untested future direction in Section 7. No source reviewed here reports measuring both correction precision on a ground-truth-preselected true-positive subset \emph{and} the net effect of the same mechanism applied unselectively to the full object population as two explicitly separated quantities — that separation is this project's specific empirical contribution to the area.

\subsection{Test-time self-correction and confirmation bias in frozen models}

The no-oracle constraint used throughout this project — correction without access to ground truth at inference time — places the entire P5 series inside the test-time-adaptation (TTA) problem class, even though the project's own working documents never use that term. Tent \citep{wang2020tent} is the canonical example of a frozen model adapted online using only its own predictions. CoTTA \citep{wang2022cotta} is more directly relevant: it diagnoses that long-horizon self-supervised adaptation on unreliable pseudo-labels causes error accumulation, and mitigates it through weight-averaging and stochastic parameter restoration — a dampening, rather than gating, strategy. This project's own ``correction cascade'' finding — objects with more than 15 self-corrections in an early no-oracle variant regress by a mean of 10.6 points, versus a mean gain of 20.6 points for objects with 15 or fewer — is an independently observed instance of the same accumulation dynamic CoTTA was built to prevent, and the eventual A2 mechanism (an object-level correction-rate circuit breaker) is, in effect, a hard-threshold version of CoTTA's soft dampening idea, arrived at without reference to the TTA literature.

\citet{arazo2020} give the general mechanism name for a self-referential-source finding observed early in this project: correction sources that are the same kind of signal as the detection signal (the model's own last-trusted mask, in several early no-oracle variants) degrade through exactly the confirmation-bias dynamic Arazo et al. describe for pseudo-labeling in classification. \citet{huang2023}, from the LLM-reasoning literature rather than vision, is included as a cross-domain structural analogue: their finding that self-correction without an external, reliable oracle frequently makes outputs worse, not better, mirrors this project's capstone finding closely. Both point to the same structural requirement — a self-correcting system's net benefit is bounded by the reliability of its trigger for ``correction needed,'' not by the quality of the correction once triggered.

\subsection{Trigger specificity under class imbalance}

This is the area where a single external source most directly explains the project's central quantitative result. \citet{saito2015} show formally that under class imbalance, a detector's strong per-instance performance does not guarantee useful net precision, because false positives on a large majority class can overwhelm true positives on a rare minority class. This project's own numbers instantiate the pattern exactly: roughly 10\% of SA-V val objects (about 30 of 293) genuinely need drift correction, and the capstone run's implied trigger false-positive rate on the remaining ``healthy'' population is approximately 68\% — high enough that, even though P5's correction is highly precise on true positives (individual gains up to 90 points J\&F), the net population-level effect is negative. This is precisely the base-rate/precision-collapse dynamic Saito and Rehmsmeier formalize, arrived at independently here through direct measurement rather than engagement with the imbalanced-classification literature. We treat this paper as the imbalanced-detection framing for the empirical results reported below.

\section{Method}

\subsection{Model, benchmark, and evaluation protocol}

All experiments use SAM2.1 Hiera-Tiny (SAM2 codebase commit \texttt{2b90b9f5ceec907a1c18\allowbreak 123530e92e794ad901a4}), frozen throughout — no weights are updated at any point in this study. A transferability check confirms \texttt{box\_centroid} does not degrade on the larger Hiera-Small checkpoint on full DAVIS val (J\&F-Mean 0.842 vs. 0.8235 for Hiera-Tiny), but the full SA-V drift-correction series was run only on Hiera-Tiny; Hiera-Small was not carried through it (see Section 6). Evaluation uses the official SA-V evaluator (\texttt{sam2/sav\_dataset/sav\_evaluator.py}), which computes \texttt{J} (region similarity), \texttt{F} (boundary accuracy), and \texttt{J\&F = (J+F)/2}, skipping the first and last annotated frame per SA-V convention. SA-V val comprises 155 videos and 293 objects; SA-V test comprises 150 videos and 278 objects. Hardware was a single NVIDIA RTX 4070 Ti (12 GB) with PyTorch 2.13.0+cu130. Every reported number comes from this official evaluator run on real predictions, never from an LLM-based or otherwise heuristic scoring proxy.

\subsection{Baseline prompting policy}

The baseline policy, \texttt{box\_centroid}, constructs a tight bounding box around the ground-truth mask on the first annotated frame plus a positive point at the mask centroid, and calls SAM2.1's \texttt{add\_new\_points\_or\_box}. It was selected from ten candidate prompting policies on a DAVIS development subset, validated on full DAVIS 2017 val (J\&F-Mean = 0.8235, official \texttt{davis2017-evaluation} evaluator, 30 sequences / 1999 frames) and then independently validated on SA-V val (J\&F = 72.1, J = 68.6, F = 75.6) and SA-V test (J\&F = 74.4). The roughly 10-point drop from DAVIS to SA-V reflects SA-V being a harder, more diverse benchmark (more objects per video, more occlusion and appearance change), not a broken pipeline.

Failure-mode analysis of the 293 SA-V val objects decomposes low-scoring cases (J\&F < 40, n = 36) into two causes: 6 of 36 (17\%) are prompt-frame failures, where the model already fails on the very frame the prompt is built from, typically because a fragmented or sparse ground-truth mask (\texttt{box\_fill\_ratio} < 0.2) produces an uninformative bounding box; 30 of 36 (83\%) are propagation-drift failures, where the model starts correctly (prompt-frame IoU $\geq$ 0.3) but degrades during propagation. An oracle experiment on these 30 propagation-drift objects — re-prompting from the ground-truth mask every fifth annotated frame — recovers J\&F from 26.1 to 84.3 (+58.2pp), establishing that the failure is a recoverable memory-persistence problem, not a perceptual ceiling.

A secondary, independent fix targets the smaller prompt-frame-collapse cause: constructing the initial box from the largest connected component of the ground-truth mask rather than its full (possibly fragmented) bounding box. This policy, referred to as P2, is a no-op for already-compact masks and was initially measured only on the 13 objects with \texttt{box\_fill\_ratio} < 0.2 (J\&F 47.1 → 61.6, +14.6pp). Section 4 reports what happens when P2 is applied to the full, unfiltered population instead.

\subsection{P5: motion-prior drift detection and correction}

Five earlier no-oracle drift-correction variants — a single self-mask reference, a 5-mask consensus reference, wider and narrower area-ratio thresholds, and SAM2's internal \texttt{object\_score\_logits} as an alternative trigger — were tried first, all re-anchoring via reinjection of the model's own recent mask content. All five converged to a narrow range, clustered within a 3.6-point band (J\&F $\in$ [23.8, 27.4] on the 30-object drift subset), far below the 84.3 oracle ceiling. The common mechanism, a ``correction cascade,'' is that the reference mask used for correction is generated by the same process that produces the error it is meant to fix, so one imprecise correction degrades the reference for the next.

P5 changes two things at once. First, detection adds a position-jump signal alongside the existing area-ratio signal: the object's centroid is extrapolated from recent frames under a constant-velocity assumption, and a frame is flagged as drift if the actual centroid deviates from the extrapolated position by more than \texttt{jump\_ratio × sqrt(EMA area)}. This is a signal the area-ratio and confidence-based triggers structurally miss: a confident identity switch onto a similarly sized distractor changes neither the mask's area nor the model's own confidence, but it does change position. Second, correction abandons reinjection of the model's own mask content entirely: on a flagged frame, P5 issues a fresh bounding-box prompt built only from the extrapolated position and the size of the last accepted mask, via \texttt{add\_new\_points\_or\_box}, letting SAM2 re-segment whatever is actually present at that location rather than continuing to trust what it was just (possibly incorrectly) tracking.

On the 30-object drift subset, untuned P5 reaches J\&F = 35.9, a 9.1-point gain over the objscore no-oracle baseline (26.8) and an 8.5-point gain over the previous best no-oracle result, consensus-of-masks (27.4). Coordinate-descent calibration of four hyperparameters (\texttt{jump\_ratio}, \texttt{low\_ratio}, \texttt{high\_ratio}, \texttt{ema\_alpha}) followed by a joint 3×3 grid check around the coordinate-descent optimum converged on \texttt{jump\_ratio = 8.0} (stable over 6.0–9.0), \texttt{low\_ratio = 0.35}, \texttt{high\_ratio = 2.5} (stable over 2.5–3.0), \texttt{ema\_alpha = 0.2}, reaching J\&F = 42.3 — the best result of the core calibration sweep, a 15.5-point gain over the objscore no-oracle baseline (a later refinement described below raises this slightly further, to 42.6). These headline 42.3/42.6 figures come from hyperparameters tuned \emph{and} evaluated on the same 30-object subset — an in-sample optimization result — so they should be read as an upper bound on this mechanism's achievable gain rather than an estimate of held-out performance; the val-to-test generalization check reported in Section 4.5 (Table 10), where this same calibrated configuration's gain falls from +15.5pp on val to +2.0pp on held-out test, is the paper's actual held-out evidence of how much this specific number should be trusted. A final parameter, \texttt{min\_area\_eps}, was confirmed at its original untuned value of 50 via a dedicated sweep. A sixth parameter, \texttt{vanish\_retry\_limit}, which caps consecutive vanish-triggered correction retries, was not part of this original four-parameter calibration; it was added later, during test-split diagnosis, and is introduced where it is evaluated (Section 4.5). On individual objects with confirmed multi-instance ambiguity, gains reach as high as +90.1pp (\texttt{sav\_022396}, roughly 15 near-identical balloons: 5.2 → 95.3).

\subsection{Mitigation families applied to P5 (and, where noted, P2)}

Applying the combined P2+P5 pipeline to the full, unselected 293-object val population — the capstone run — produced the paper's central diagnostic result (Section 4.3), motivating twelve further attempts, organized into four families.

\textbf{Frame-level gating.} A1 requires both the area-ratio and position-jump signals to fire simultaneously (an AND-gate) rather than either alone (the original OR-gate). B1 checks the \emph{result} of a completed correction against recent size/position history and, if implausible, substitutes a conservative ``hold position'' fallback box (no velocity extrapolation) instead of accepting the correction.

\textbf{Object-level circuit breakers.} A2 tracks each object's cumulative correction rate (\texttt{n\_corrections / frames\_elapsed}) and, after a 3-correction warmup, permanently disables further correction for that object once the rate exceeds a threshold (0.15, tuned on the 30-object subset and confirmed on full val). D1 re-sweeps this threshold (0.10, 0.20, 0.25) directly on full val. D2 lowers the warmup-corrections requirement (1 or 2 instead of 3). D3 replaces the cumulative rate with a rate computed over a sliding window of recent frames. D4 stacks A2 and A1 together.

\textbf{True pre-selection (E1).} Rather than reacting after corrections have already begun (as A2 does), E1 computes risk features from an uncorrected dry-run propagation pass, before any correction is attempted: (a) raw trigger rate and area/position-jump amplitude statistics; (b) the fraction of an object's total frames elapsed before its first trigger (\texttt{first\_trigger\_frac}); (c) the fraction of an object's triggers occurring in its first 10\% or 20\% of frames. All three were evaluated by correlating dry-run features against the capstone-minus-P2-only per-object delta on the same 293 val objects.

\textbf{Pre-selection from an intra-frame introspective signal, DAM4SAM-inspired.} E1's three feature families all summarize an object's \emph{trajectory} over time. A structurally different candidate signal comes from a single frame in isolation: SAM2 computes three candidate masks and their IoU estimates at every tracking step but keeps only the best-scoring one, discarding the rest before any caller can see them. Following the introspective distractor-detection idea in DAM4SAM \citep{videnovic2024} — which uses disagreement between a tracker's primary and alternative candidate masks as a distractor signal — the discarded candidates were captured non-invasively, by wrapping the tracker's own mask-selection step at runtime rather than modifying the underlying model, and used to derive six per-object features from a dedicated uncorrected dry-run pass over all 293 val objects: mean and maximum bounding-box-area divergence between the primary and best-alternative mask, the fraction of frames with divergence above a fixed threshold, the frame-position of an object's first divergence event, and the mean and minimum gap between the primary and second-best mask's IoU estimates. These were evaluated by the same correlation methodology as E1, against the same capstone-minus-P2-only delta on the same 293 objects.

\textbf{Correction softening (F1).} Rather than gating whether or when correction happens, F1 changes \emph{how} a correction is applied: the reprompt box's center is blended between the fully extrapolated position and the model's own (possibly drifted) current centroid, \texttt{(1-w)·predicted\_centroid + w·raw\_centroid}, instead of a full hard override, at blend weights \texttt{w} swept from 0.10 to 1.00.

\subsection{Evaluation caveats}

Every number reported in Section 4 comes from the official SA-V evaluator on real predictions. No formal significance testing — bootstrap confidence intervals, multiple seeds, or p-values — was computed anywhere in this series; conclusions rest on the consistency of effect direction across multiple independent manipulations and, where both splits were run, agreement between val and test. This is a deliberate deviation from the project's own stated evaluation protocol, which specifies bootstrap 95\% confidence intervals and three-seed repetition as a precondition for publishing a claimed improvement; we discuss this explicitly as a limitation in Section 6 rather than retrofitting statistics that were not computed.

\section{Results}

\subsection{Baseline and failure decomposition}

\begin{table}[!htbp]
\centering
\caption{Baseline transfer from DAVIS to SA-V.}
\label{tab:1}
\small
\begin{tabularx}{\textwidth}{@{}LLRrrr@{}}
\toprule
Setting & Split & n & J\&F & J & F \\
\midrule
\texttt{box\_centroid}, Hiera-Tiny & DAVIS 2017 val (official) & 30 seq / 1999 frames & 82.4 (0.8235) & 78.1 & 86.6 \\
\texttt{box\_centroid}, Hiera-Small & DAVIS 2017 val (official) & 30 seq / 1999 frames & 84.2 (0.842) & 80.1 & 88.3 \\
\texttt{box\_centroid}, Hiera-Tiny & SA-V val (official) & 155 videos / 293 objects & 72.1 & 68.6 & 75.6 \\
\texttt{box\_centroid}, Hiera-Tiny & SA-V test (official) & 150 videos / 278 objects & 74.4 & --- & --- \\
\bottomrule
\end{tabularx}
\end{table}

\begin{table}[!htbp]
\centering
\caption{SA-V val failure decomposition (n = 293 objects).}
\label{tab:2}
\small
\begin{tabularx}{\textwidth}{@{}Lrr@{}}
\toprule
Category & Fraction & Mean J\&F \\
\midrule
\texttt{box\_fill\_ratio} < 0.2 (sparse/fragmented mask) & n = 13 & 47.1 \\
\texttt{box\_fill\_ratio} 0.2–0.5 & n = 72 & 71.1 \\
\texttt{box\_fill\_ratio} $\geq$ 0.5 (compact) & n = 208 & 74.0 \\
Prompt-frame failures (among J\&F < 40, n = 36) & 6/36 (17\%) & --- \\
Propagation-drift failures (among J\&F < 40, n = 36) & 30/36 (83\%) & --- \\
\bottomrule
\end{tabularx}
\end{table}

\subsection{Oracle ceiling and the no-oracle self-referential ceiling}

\begin{table}[!htbp]
\centering
\caption{Recoverability of the 30-object propagation-drift subset.}
\label{tab:3}
\small
\begin{tabularx}{\textwidth}{@{}LLLr@{}}
\toprule
Variant & Detection signal & Correction source & J\&F \\
\midrule
Baseline (\texttt{box\_centroid}, single prompt) & --- & --- & 26.1 \\
Oracle periodic re-prompt (every 5th frame) & --- & ground-truth mask & \textbf{84.3 (+58.2pp)} \\
Single-mask self-reference & area-ratio vs. EMA (0.3–3.0) & model's own last mask & 25.8 \\
Consensus-5 self-reference & area-ratio vs. EMA (0.3–3.0) & majority vote of last 5 masks & 27.4 \\
Wide threshold & area-ratio vs. EMA (0.15–5.0) & model's own last mask & 23.8 \\
Tight threshold & area-ratio vs. EMA (0.5–2.0) & model's own last mask & 26.7 \\
\texttt{object\_score\_logits} trigger & internal presence score (threshold 0.0) & model's own last mask & 26.8 \\
\bottomrule
\end{tabularx}
\end{table}

All five no-oracle self-referential variants cluster tightly, within a 3.6-point band (J\&F $\in$ [23.8, 27.4]), and stay far below the 84.3 oracle ceiling. Stratifying the single-mask variant by number of self-corrections shows a bimodal ``correction cascade'': objects with $\leq$15 corrections gain a mean +20.6pp (n=10), objects with >15 corrections lose a mean $-$10.6pp (n=20).

An external-correction pilot (n = 6 objects, Claude as reviewer, no API integration, no ground truth) tested whether an independent correction source alone helps: objscore-only baseline J\&F = 18.5 fell to 10.2 ($-$8.3pp) once externally drawn boxes replaced self-mask correction, including a 48.8-point collapse on a well-tracked control object (\texttt{sav\_010291}, 84.6 → 35.8). Four of six objects were correctly flagged by the reviewer as genuinely ambiguous (multi-instance identity confusion) and abstained rather than guessed.

\subsection{P5: calibration, and the curated-subset-versus-full-population contrast}

\begin{table}[!htbp]
\centering
\caption{P5 calibration path on the 30-object drift subset.}
\label{tab:4}
\small
\begin{tabularx}{\textwidth}{@{}LrR@{}}
\toprule
Configuration & J\&F & $\Delta$ vs. best prior no-oracle (27.4) \\
\midrule
Untuned P5 (area-ratio + position-jump, default params) & 35.9 & +8.5pp \\
P5, \texttt{jump\_ratio = 8.0} only calibrated & 39.3 & +11.9pp \\
P5, fully calibrated (\texttt{jump\_ratio=8.0, low\_ratio=0.35, high\_ratio=2.5, ema\_alpha=0.2}) & \textbf{42.3} & \textbf{+14.9pp} \\
\bottomrule
\end{tabularx}
\end{table}

\emph{Table 4's $\Delta$ column is computed against 27.4 (the consensus-5 variant), the single best-performing configuration among the five prior no-oracle attempts shown in Table 3. The narrative baseline used for the headline ``+15.5pp'' figure elsewhere in this section and in Table 5 is instead 26.8 (the objscore variant): both are real reference points from Table 3, but 27.4 is the best of the prior attempts overall, while 26.8 is the specific configuration P5 finally replaces in the deployed, capstone comparison.}

Relative to the objscore no-oracle baseline (26.8), fully calibrated P5 gains +15.5pp. Per-object, 21/30 objects improve, 4/30 regress (worst: $-$4.3pp, no catastrophes), and 5/30 are unchanged. The largest gains concentrate exactly on objects with independently confirmed multi-instance identity ambiguity: \texttt{sav\_022396} (5.2 → 95.3, +90.1pp), \texttt{sav\_035221/001} (11.4 → 63.3, +51.9pp), \texttt{sav\_047467/002} (24.2 → 64.1, +39.9pp).

\begin{table}[!htbp]
\centering
\caption{(Headline) Curated true-positive subset versus full, unselected population — the paper's central empirical result.}
\label{tab:5}
\small
\begin{tabularx}{\textwidth}{@{}LrLrL@{}}
\toprule
Population & n & Condition & J\&F & $\Delta$ vs. no-correction reference \\
\midrule
GT-preselected genuine drift-failure objects & 30 & Fully calibrated P5 & 42.3 & \textbf{+15.5pp} vs. objscore no-oracle baseline (26.8) \\
Full, unselected SA-V val & 293 & P2+P5 combined, no gating (capstone) & 68.7 & \textbf{$-$3.4pp} vs. plain \texttt{box\_centroid} baseline (72.1) \\
\bottomrule
\end{tabularx}
\end{table}

Per-object outcomes on the full population: 44/293 improve, 178/293 regress, 71/293 unchanged. The improvements are the same objects P5 was calibrated for (\texttt{sav\_022396}: 5.0 → 95.3, +90.3pp; \texttt{sav\_022456}: 0.4 → 91.2, +90.8pp; \texttt{sav\_045706}: 5.6 → 95.3, +89.7pp), but the regressions are more numerous and comparably severe, striking objects the baseline already tracked well: \texttt{sav\_014969/000} (89.6 → 1.2, $-$88.4pp), \texttt{sav\_024930/001} (88.2 → 6.2, $-$82.0pp), \texttt{sav\_021068/000} (78.3 → 10.7, $-$67.6pp), and at least 15 further objects lose more than 35 points. Roughly 10\% of val objects (about 30 of 293) genuinely need drift correction; the trigger's implied false-positive rate on the remaining ``healthy'' $\sim$90\% is approximately 68\%.

Two independent diagnostics corroborate that the trigger, not the correction step, is the problem. First, 200/293 objects (68\% of the whole population) fire the position-jump trigger at least once, and this subgroup has a mean delta of $-$5.17pp — worse than the population average — indicating the position signal is not, in practice, more specific to genuine identity switches than the area-ratio signal, despite the design intuition that a purely spatial signal should better resist confusion between visually identical instances. Second, correction rate per object is a cleanly monotonic predictor of outcome:

\begin{table}[!htbp]
\centering
\caption{Correction rate versus outcome, full SA-V val (n = 293). Per-bucket improved/regressed sub-counts as logged during the retrospective analysis are omitted here because they do not sum to the same population-wide improved/regressed totals reported for the capstone run above (a discrepancy present in the underlying experimental log that this paper does not attempt to silently resolve); the mean-delta trend below is unambiguous and is what motivated A2's design.}
\label{tab:6}
\small
\begin{tabularx}{\textwidth}{@{}Lrr@{}}
\toprule
Correction rate (\texttt{n\_corrections / frames\_elapsed}) & n objects & Mean $\Delta$ \\
\midrule
< 0.05 & 88 & +3.28pp \\
0.05–0.15 & 88 & $-$3.26pp \\
0.15–0.30 & 67 & $-$5.70pp \\
> 0.30 & 50 & $-$11.87pp \\
\bottomrule
\end{tabularx}
\end{table}

\subsection{P2 isolated on the full, unselected population — a second illustration of the same pattern}

\begin{table}[!htbp]
\centering
\caption{P2 (largest-connected-component prompting), curated subset versus full population.}
\label{tab:7}
\small
\begin{tabularx}{\textwidth}{@{}LrrL@{}}
\toprule
Population & n & J\&F & $\Delta$ vs. \texttt{box\_centroid} \\
\midrule
Curated: sparse-mask objects only (\texttt{box\_fill\_ratio} < 0.2) & 13 & 61.6 & \textbf{+14.6pp} (4/13 substantially improved up to +90.8pp, 8/13 unchanged, 1/13 regressed $-$11.8pp) \\
Full SA-V val (isolated, no P5) & 293 & 72.5 & +0.4pp \\
Full SA-V test (isolated, no P5) & 278 & 74.1 & \textbf{$-$0.3pp} \\
\bottomrule
\end{tabularx}
\end{table}

On the curated subset, P2 looks like a free, universally safe fix. On the full, unselected population it is high-variance and not risk-free: val shows 21/293 objects improving (up to +90.8pp) against 24/293 regressing (up to $-$48.5pp, \texttt{sav\_024172/003}: 72.5 → 24.0); test shows 16/278 improving (up to +50.7pp) against 31/278 regressing (up to $-$38.2pp, \texttt{sav\_042227/000}: 86.4 → 48.2), with the regression tail dominating on test (net $-$0.3pp). P2 on the 13-object test subset used for the original curated measurement is byte-identical to baseline (no fragmented masks present), a 0.0pp isolated effect fully inert on that split, in contrast to +14.6pp on the analogous val subset. Both ``confirmed'' positive results in this project's series (P2 on its curated subset, P5 on its curated subset) substantially or fully lose their effect once measured on a full, unselected population.

\subsection{Mitigation attempts}

\begin{table}[!htbp]
\centering
\footnotesize
\caption{Frame-level and object-level mitigations, full SA-V val and (where run) test.}
\label{tab:8}
\footnotesize
\begin{tabularx}{\textwidth}{@{}>{\raggedright\arraybackslash\hsize=1.3\hsize}X>{\raggedright\arraybackslash\hsize=1.0\hsize}Xr>{\raggedleft\arraybackslash\hsize=1.0\hsize}X>{\raggedleft\arraybackslash\hsize=0.7\hsize}Xrr@{}}
\toprule
Configuration & Family & val J\&F & val $\Delta$ & val improved/ regressed & test J\&F & test $\Delta$ \\
\midrule
Baseline (\texttt{box\_centroid}) & --- & 72.1 & --- & --- & 74.4 & --- \\
P2 isolated & prompt-only & 72.5 & +0.4pp & --- & 74.1 & $-$0.3pp \\
Capstone OR-gate (P2+P5, no gate) & reference & 68.7 & $-$3.4pp & 44/178 & not run & --- \\
A1: AND-gate (both triggers required) & frame-level gating & 69.0 & $-$3.1pp & 50/167 & not run & --- \\
B1: plausibility-check on correction result & frame-level gating & 68.2 & $-$3.9pp & 48/199 & not run & --- \\
A2: rate circuit breaker (threshold 0.15) & object-level breaker & 71.1 & $-$1.0pp & 49/189 & 72.2 & $-$2.2pp \\
D1: rate threshold 0.10 & object-level breaker & 71.2 & $-$0.9pp & --- & not run & --- \\
D1: rate threshold 0.20 & object-level breaker & 70.5 & $-$1.6pp & --- & not run & --- \\
D1: rate threshold 0.25 & object-level breaker & 70.3 & $-$1.8pp & --- & not run & --- \\
D2: warmup=1 (on D4 config) & object-level breaker & 70.9 & $-$1.2pp ($-$0.4pp vs. D4) & --- & not run & --- \\
D4: A2 (0.15) + A1 stacked & object-level + frame-level & \textbf{71.3} & \textbf{$-$0.8pp} & --- & \textbf{72.6} & \textbf{$-$1.8pp} \\
F1: soft-nudge, w=0.10 (on capstone, no A1/A2) & correction softening & 67.9 & $-$4.2pp & --- & not run & --- \\
\bottomrule
\end{tabularx}
\end{table}

D3 (sliding correction-rate window), E1 (all three pre-selection variants), and the DAM4SAM-inspired multi-mask divergence check never reached a full corrected run because cheap pilots or correlations on already-collected or purpose-run dry-run data were unambiguously negative (D3) or showed no usable signal (E1, multi-mask divergence); Table 9 reports those checks.

\begin{table}[!htbp]
\centering
\footnotesize
\caption{Cheap-gate checks that closed a direction without a full run.}
\label{tab:9}
\small
\begin{tabularx}{\textwidth}{@{}LLL@{}}
\toprule
Direction & Check & Result \\
\midrule
D3: sliding correction-rate window (sizes 20/30/50/100) & n=30 pilot vs. cumulative A2 (40.2) & Best window (100) still $-$0.4pp; no window beat cumulative; closed without full run \\
E1-rate: 6 dry-run features (trigger rate, amplitude) & Multiple regression, all 293 val objects, vs. capstone-minus-P2-only delta & R$^2$ = 0.026; single-feature buckets non-monotonic; no usable signal \\
E1-timing: \texttt{first\_trigger\_frac} as live gate & Exact file-copy reconstruction of gated pipeline, 293 val objects, thresholds T $\in$ [0.02, 0.70] & Curve saturates exactly at the already-known, non-generalizing P2-only ceiling (72.5), never exceeds it \\
E1-timing: fraction of triggers in first 10\%/20\% of frames & Correlation, same 293 objects & r = $-$0.111 / $-$0.109 (all), non-monotonic buckets; combined 9-feature R$^2$ = 0.044 \\
Multi-mask divergence (DAM4SAM-inspired): 6 intra-frame introspective features & Multiple regression, dedicated dry-run pass, all 293 val objects, vs. capstone-minus-P2-only delta & R$^2$ = 0.028 (materially unchanged from E1's 0.026); strongest single feature r = $-$0.131 (p = 0.024) does not survive a Spearman rank test ($\rho$ = $+$0.054, p = 0.35, sign flips) and is 1 of 6 features tested without multiple-comparisons correction; no usable signal \\
\bottomrule
\end{tabularx}
\end{table}

F1 (soft-nudge) was the only mitigation whose n=30 pilot was steeply negative at every tested weight yet was still run at full scale, on the reasoning that a mechanism changing \emph{how} correction is applied (rather than whether or when) could plausibly help the 263 objects outside the curated subset even while hurting the 30 inside it. It did not: full-val J\&F = 67.9 (w=0.10), worse than the unmitigated OR-gate capstone (68.7) and the single worst result of the entire series.

\begin{table}[!htbp]
\centering
\caption{P5 val-to-test generalization gap on the curated drift subset.}
\label{tab:10}
\small
\begin{tabularx}{\textwidth}{@{}lrrrrL@{}}
\toprule
Split & n & Baseline J\&F & P5 J\&F & $\Delta$ & Improved/Regressed/Unchanged \\
\midrule
val & 30 & 26.8 & 42.3 & +15.5pp & 21/4/5 \\
test & 25 & 26.9 & 28.9 & +2.0pp & 10/10/5 \\
\bottomrule
\end{tabularx}
\end{table}

On test, P5 (with hyperparameters calibrated only on val) produces two catastrophic regressions absent on val: \texttt{sav\_017171/000} (87.5 → 53.7, $-$33.8pp, driven by the target rapidly growing in size while the reprompt box retains a stale, undersized reference) and \texttt{sav\_004755/000} (34.3 → 0.0, $-$34.3pp, driven by the target vanishing and P5 confidently re-anchoring on the wrong object, which then self-reinforces across all 68 remaining annotated frames). Both catastrophes originate from the area-ratio trigger (\texttt{position\_triggers = 0} in both cases). A targeted fix limiting consecutive vanish-triggered retries (\texttt{vanish\_retry\_limit}) improves val slightly (42.3 → 42.6 at limit 10) but leaves both test catastrophes completely unchanged, because each stems from a single confidently wrong first correction rather than a chain of retries.

\subsection{Final result}

\begin{table}[!htbp]
\centering
\caption{Final comparison: best mitigation versus baseline, both splits.}
\label{tab:11}
\begin{tabular}{@{}lrrr@{}}
\toprule
Split & Baseline & D4 (best mitigation) & $\Delta$ \\
\midrule
val & 72.1 & 71.3 & $-$0.8pp \\
test & 74.4 & 72.6 & $-$1.8pp \\
\bottomrule
\end{tabular}
\end{table}

No configuration across the full series — the five early no-oracle variants (Table 3), P2, P5, the capstone, A1, B1, A2, D1–D4, three E1 variants, the DAM4SAM-inspired multi-mask divergence check, and F1 — exceeds the plain \texttt{box\_centroid} baseline on either split.

D4 is reported as the single best mitigation for concreteness, but its margin over the next-closest configurations is narrow: on val, D4 ($-$0.8pp), D1 at threshold 0.10 ($-$0.9pp), and A2 ($-$1.0pp) sit within 0.1–0.2 points of one another (Table 8) — a gap plausibly within measurement noise on a 293-object population for which no confidence intervals, multiple seeds, or significance tests were computed (Section 3.5). The robust claim is the qualitative one: D4, A2, and D1@0.10 all remain clearly below the uncorrected baseline; their exact relative ordering is not something this evaluation protocol can support.

\section{Discussion}

The unifying lens for every result in Section 4 is the distinction between correction precision (how well a correction matches ground truth when it is genuinely needed) and trigger specificity (how rarely the mechanism fires on objects that did not need it). Every attempt before the capstone run measured only the first quantity, on a population already known to need correction; the capstone run was the first point at which the second quantity was measured at all, and it showed that specificity, not precision, is what determines whether the mechanism is safe to deploy. P5's correction step is excellent by the only standard the earlier, curated-subset experiments could apply — it beats every prior no-oracle correction source by a wide margin and approaches oracle-level gains on genuine multi-instance identity-switch cases. None of that predicts what happens when the same trigger logic is exposed to a population where roughly 90\% of objects do not need correction at all: a trigger with even a moderate false-positive rate on that majority produces more total harm than the minority's gains can offset, which is exactly the arithmetic \citet{saito2015} formalize for imbalanced classification and exactly what the 178-versus-44 object-level split shows.

This also explains why object-level mechanisms (A2, D4) outperformed frame-level ones (A1, B1) even though none of them closed the gap to baseline. A1's AND-gate assumes the area-ratio and position-jump signals are at least partially independent evidence of drift, so requiring both should suppress false positives more than either alone. The data reject that assumption: 68\% of the whole val population fires the position-jump trigger at least once, and this subgroup's mean outcome ($-$5.17pp) is worse than the population average, because both signals respond to the same underlying legitimate events — a fast scale change or fast motion moves centroid and area simultaneously. Requiring agreement between two correlated, jointly non-specific signals recovers only 0.3pp over the OR-gate. B1's plausibility check on the correction's \emph{result} fails for a related but distinct reason: its ``safe'' fallback (hold position, no velocity extrapolation) is only safe for slow or stationary objects, and the population where the trigger fires most is disproportionately fast-moving objects, precisely where a non-extrapolating fallback lags the true position most. Because velocity resets to zero after any correction including a rejected one, this fallback creates a new position jump on the very next frame, triggering another correction and another rejection — a self-reinforcing mis-anchoring loop that produced more regressions (199/293) than either gate it was meant to improve on. A2 and D4 instead act on an object's accumulated correction-rate history rather than any single frame's trigger or result. This is a genuinely different kind of signal: it does not attempt to distinguish a legitimate fast-motion event from a real drift event at the moment either occurs (both A1 and B1's failure mode), but instead recognizes, after the fact, that an object generating corrections at a chronically high rate is behaving like the population where P5 causes net harm, and stops correcting that object regardless of the reason for any individual trigger. This does not reduce the trigger's false-positive \emph{rate}; it caps the \emph{damage} a chronically triggering object can accumulate, which is why the worst single regression under A2 ($-$65.7pp) is markedly smaller than under the OR-gate ($-$88.4pp) while every major true-positive win is preserved untouched — short, accurate corrections never accumulate enough rate to trip the breaker.

E1's pre-selection attempts found essentially no signal (R$^2$ = 0.026 for raw rate/amplitude features, non-monotonic buckets throughout) for the same underlying reason A1 failed: the features available from an uncorrected dry-run pass — how often and how severely a threshold is crossed — measure the same thing the live trigger already measures, just without acting on it. They cannot distinguish ``the trigger is correctly detecting an identity switch'' from ``the trigger is incorrectly reacting to ordinary fast motion or scale change,'' because that distinction is exactly the discrimination problem the trigger itself fails to solve; recomputing the same statistic on the same uncorrected trajectory does not introduce new information capable of resolving it. The one E1 variant that did show a real, non-zero correlation — the timing of an object's first trigger — turned out, on reconstruction across the full threshold range, to saturate exactly at the already-known, non-generalizing P2-only ceiling (72.5) and never exceed it: the gate was not discovering a new source of signal, it was progressively disabling P5 back toward P2-only as the threshold tightened, and P2-only is already known to reverse sign on test.

A later check pursued the same question with a categorically different signal — not the trajectory of the selected mask over time, but SAM2's own multi-hypothesis disagreement within a single frame (the divergence between its primary and best-alternative candidate mask, in the spirit of DAM4SAM's introspective distractor detection) — and found the same absence of signal (R$^2$ = 0.028, materially unchanged from E1's 0.026; its one nominally significant single feature does not survive a rank test and is uncorrected for multiple comparisons). This is some evidence, though not proof given only two feature families were tried, that the discrimination this project needs is not recoverable from cheap, no-oracle introspective signals of either general kind: neither a mask's history across frames nor the tracker's own competing hypotheses within a frame separates a genuine identity switch from an ordinary motion or scale event that happens to cross the same trigger thresholds.

F1's failure is the most mechanistically direct of the twelve attempts. P5's entire advantage over the earlier self-referential correction variants comes from refusing to trust the model's own current, possibly drifted position and instead re-segmenting from a purely extrapolated one. Blending any fraction of the model's own centroid back into the reprompt box directly reintroduces the untrustworthy signal that P5 was built to discard, and it does so at every correction, including the true positives that account for essentially all of P5's benefit. That the gentlest tested blend (w = 0.10) already cost 5.6 points on the curated true-positive subset, with no plateau near zero, is consistent with this: there was no safe amount of reintroducing the compromised position, because the position is precisely what the mechanism must not use.

A second theme cuts across several results independently: measurements on a curated or GT-preselected subset do not reliably predict full-population, unselected effect, and this held in both directions and on both mechanisms tried this way. P5's val gain (+15.5pp on the true-positive subset) fell to +2.0pp on the analogous held-out test subset, with two catastrophic regressions on test absent on val. P2's isolated effect flipped sign between splits (+0.4pp val, $-$0.3pp test). And within a single split, the D2 warm-up-threshold pilot on the 30-object subset showed an essentially free gain (+0.2pp) that reversed to a $-$0.4pp regression at full scale, because the 30-object subset — selected for having multiple sequential corrections — structurally excludes the scenario (a single needed correction, blocked by one earlier accidental trigger) that caused the full-population regression. None of these reversals required a different mechanism to produce; the same mechanism, measured on a differently composed population, gave a different-signed answer each time.

A final scoping note applies to the paper's central claim. The specific numbers reported here — a $\sim$10\% positive-class base rate, a $\sim$68\% trigger false-positive rate, and a $-$3.4pp net effect — are properties of this particular mechanism (P5's area-ratio and position-jump triggers), this particular frozen checkpoint (SAM2.1 Hiera-Tiny), and this particular benchmark (SA-V), and should not be assumed to transfer numerically to a different correction mechanism, model, or dataset. The qualitative pattern is a separate matter: any no-oracle correction mechanism whose trigger is validated for precision only on a population already known to need correction, and never measured for its false-positive behavior on the (typically much larger) population that does not, is structurally exposed to the same base-rate arithmetic \citet{saito2015} formalize for imbalanced classification generally — the mechanism-, model-, and dataset-specific numbers will differ, but the underlying requirement to measure specificity on an unselected population before deployment does not depend on any of those specifics. We therefore expect the trigger-specificity bottleneck itself, not its exact magnitude, to generalize beyond this study's setup. A direct methodological consequence follows: any no-oracle correction mechanism should report both its precision on a curated, ground-truth-preselected true-positive subset \emph{and} its net effect when applied unselectively to the full target population before being considered validated for deployment — this paper's own P5 and P2 results would each have looked unambiguously positive had only the first measurement been taken.

\section{Limitations}

\textbf{Model scale.} The full drift-correction series (P5 and all twelve mitigations) was run only on SAM2.1 Hiera-Tiny. Hiera-Small was checked only for baseline \texttt{box\_centroid} transferability on DAVIS val (J\&F-Mean 0.842 vs. 0.8235), not carried through any part of the SA-V drift-correction investigation. Whether the trigger-specificity problem this paper diagnoses is checkpoint-size-dependent is unknown.

\textbf{Single-benchmark generalization.} All drift-correction results are measured on SA-V. Whether the same false-positive dynamics hold on other VOS benchmarks with different object density, occlusion frequency, or motion statistics is untested.

\textbf{Uneven exploration depth across mitigation attempts.} The twelve mitigation attempts did not all receive comparably thorough exploration. P5 itself, and the A1/A2/D1/D4 family that became the eventual best result, received multi-parameter coordinate-descent or grid-search calibration (Section 3.3; Table 8). D2, D3, F1, E1, and the DAM4SAM-inspired multi-mask divergence check instead relied on a cheap n=30-pilot gate or a correlation on collected dry-run data before any full-scale corrected commitment; D3, all three E1 variants, and the multi-mask divergence check never received a full-scale \emph{corrected} run at all once their pilot or correlation checks came back negative (Table 9), and even where a full-scale corrected run did happen, only a single configuration was tested for D2 (warmup=1) and for F1 (w=0.10) — this despite F1's own n=30 sweep being non-monotonic in the blend weight (weight=0.75 scored worse, J\&F=32.6, than weight=1.00, J\&F=35.2, not a simple monotonic worsening as weight increased). The paper's ``twelve mitigation attempts'' framing should not, therefore, be read as twelve comparably rigorous investigations.

\textbf{Small-sample-pilot protocol has a documented failure mode.} The project's standard practice — a cheap n=30 pilot or correlation on existing or dedicated dry-run data before committing to an expensive (5–9 hour) full-dataset run — correctly predicted direction for most mitigations (D1, D3, all three E1 variants, the multi-mask divergence check, F1) but failed at least once: D2's n=30 pilot showed an essentially cost-free threshold change that reversed to a real regression at full scale, because the curated 30-object subset does not contain the failure scenario (a single accidental early trigger permanently blocking a later needed correction) that caused the full-population loss. Readers should not treat an n=30 pilot result, on its own, as conclusive for this class of mechanism.

\textbf{No formal statistical significance testing.} No bootstrap confidence intervals, multiple-seed repetition, or p-values were computed anywhere in this series, for any comparison. All reported deltas are exact point estimates from single runs against the official evaluator. Conclusions in this paper rest on consistency of effect direction across independent manipulations and, where both splits were measured, agreement between val and test — not on formal significance. This is a genuine limitation and is stated here rather than retrofitted with invented statistics.

\section{Future Work}

The following are literature-derived candidate directions that target the diagnosed bottleneck — trigger specificity, not correction precision — and were \textbf{not} implemented or tested by this project. They are proposals only, drawn from the sources reviewed in Section 2, and should not be read as extensions of this paper's own results.

\begin{enumerate}
\item \textbf{Object-level introspective memory gating}, in the style of DAM4SAM \citep{videnovic2024}: rather than gating on a trigger signal computed \emph{after} a correction attempt (as B1 did) or on an accumulated correction rate (as A2 did), decide whether to trust and write a given frame's propagated mask into memory \emph{before} it can contaminate future frames — potentially preventing the single-bad-frame failure documented for \texttt{sav\_004755}, where one incorrect correction became a permanent, self-reinforcing anchor. This project tested DAM4SAM's underlying introspective signal — the divergence between SAM2's primary and best-alternative candidate mask — in a narrower role, as a static per-object pre-selection feature, and found no usable signal (Section 4.5). That is a different question from the one posed here: DAM4SAM's actual gate acts continuously, frame by frame, on what enters a persistent memory bank, not as a single up-front decision of whether to run correction on an object at all. A later, separate follow-up outside this paper's twelve mitigation attempts has since tested exactly that dynamic, memory-level role, applied to plain, uncorrected propagation with no reprompt correction at all: skipping the write of a suspected-distractor frame into SAM2's non-conditioning memory bank. A curated n=30 known-drift pilot cleared its pre-committed gate by a wide margin (threshold 0.2, J\&F=39.4 against a 26.1 baseline), but the same threshold produced a net regression on the full 293-object validation population (J\&F=71.3 against a 72.1 baseline, 0.8 points below) — the same curated-subset-versus-full-population divergence this paper documents for P5. The more ambitious half of DAM4SAM's mechanism that this simpler skip-write gate cannot reach — injecting a persistent distractor-resolving anchor frame into SAM2's memory-attention window outside its normal recency schedule — remains only partly tested: the original concern that SAM2's learned positional-encoding table has no entry for an arbitrary out-of-schedule distance has since been addressed at the design stage (a proposed route via SAM2's separate, continuously-encoded object-pointer channel, a code path the frozen model already exercises on out-of-range distances for its own conditioning-frame pointer on long videos). A narrower variant has since been implemented and piloted: when the previous frame is suspected of a distractor, the content of a dynamically updated anchor object pointer, taken from the last confirmed-clean frame, replaces the oldest slot (\texttt{t\_diff} = 15) of SAM2's object-pointer window, leaving the total number of pointer tokens unchanged. On the n=30 known-drift subset it reached J\&F=27.8 against the baseline's 26.1 (5 objects improved, 2 regressed, 23 unchanged), short of its pre-committed pilot gate of 29.1, so it was closed without a full-validation run. That variant keeps the slot's fixed positional distance; injection at the anchor frame's true, potentially much larger distance — the variant the continuous-encoding argument above actually supports — has not been tested, and this negative pilot says nothing about it.
\item \textbf{Multi-path propagation}, in the style of SAM2Long \citep{ding2024}: maintain several candidate propagation pathways per object continuously, pruned by cumulative confidence, instead of this project's single-hypothesis, detect-then-hard-reprompt design, structurally avoiding irreversible commitment to one bad correction.
\item \textbf{Learned, appearance-based distractor discrimination combined with P5's positional signal}, in the style of DaSiamRPN \citep{zhu2018}: P5 uses position specifically because appearance cannot separate visually identical instances; a hybrid signal — a positional prior gated or weighted by a learned distractor-awareness score — was not tried and could plausibly raise specificity on the healthy population without losing sensitivity on genuine multi-instance cases.
\item \textbf{Object-level re-detection with multi-hypothesis dynamic programming}, in the style of Siam R-CNN \citep{voigtlaender2020}: track candidate distractor identities alongside the target through time and resolve ambiguity retrospectively via a tracklet-association step, rather than a per-trigger, forward-only, single-decision reprompt.
\item \textbf{Object-level memory reading instead of box-level correction}, in the style of Cutie \citep{cheng2023cutie}: Cutie's object-query-based memory reading is a categorically different mechanism (top-down object representation versus bottom-up pixel/box matching) not accessible to this project's frozen, prompt-only intervention surface, but it could in principle be combined with a SAM2-based pipeline as an external re-identification module.
\item \textbf{Explicit precision-recall-based threshold calibration against the known $\sim$10\% positive-class base rate}, following the framing in \citet{saito2015}: this project's trigger thresholds were tuned by maximizing J\&F on a GT-preselected true-positive subset — implicitly a 100\% base rate — never against a precision/recall operating point computed on the true $\sim$10\% base rate of the full population. A threshold explicitly chosen to hold trigger precision above a target level on the unselected population was not tried.
\item \textbf{Long-term, tiered memory consolidation as an alternative to reprompting}, in the style of XMem \citep{cheng2022xmem}: rather than detecting drift and re-anchoring via a fresh box prompt, maintain a consolidated long-term memory bank inherently more resistant to single-frame corruption, addressing the root cause (buffer rollover losing the true anchor) rather than reacting to its symptom — though this requires architectural access this project's frozen-checkpoint design does not use.
\end{enumerate}

\section{Conclusion}

A frozen SAM2.1 Hiera-Tiny, prompted well and evaluated on the official SA-V benchmark, scores J\&F = 72.1 on val and 74.4 on test with no correction at all. A motion-prior drift-correction mechanism designed to close a documented 58.2-point oracle-recoverable gap achieves a 15.5-point gain on the objects it was built for, but is net-harmful — 3.4 points below baseline — when applied without selection to the full, unselected population it would actually need to run on in deployment. Twelve subsequent mitigation attempts, spanning frame-level gating, object-level circuit breakers, feature-based pre-selection, and correction softening, narrow this gap without closing it; the best result across the entire series remains 0.8 points below baseline on val and 1.8 points below on test. That best configuration, D4, is not decisively better than several close alternatives — A2 ($-$1.0pp) and D1@0.10 ($-$0.9pp) land within 0.1–0.2pp of it on val, a margin plausibly within measurement noise given the absence of any computed confidence intervals (Section 3.5) — so the finding that matters is that no configuration in this family closes the gap to baseline, not the precise ranking among the closest few. We read this not as a failure to find the right hyperparameter but as evidence that the project's own framing is correct: trigger specificity under a roughly 10\% positive-class base rate, not correction precision on the objects that need it, is the binding constraint on whether a no-oracle drift-correction mechanism is safe to deploy on a frozen video segmentation model. The practical recommendation from this series is to deploy \texttt{box\_centroid} without drift correction; none of the investigated modifications justifies the added inference cost and complexity of a fresh-box reprompt mechanism given the measured held-out results. More broadly, we recommend that any no-oracle correction mechanism report both curated-subset precision and full-population net effect before being considered validated for deployment. We report the full series, including every closed-negative direction, so that the specific things that did not work — and, more importantly, the reasons they did not — are available to anyone attempting the untested directions in Section 7.

\section*{Data Availability Statement}

SA-V (Segment Anything Video) is a dataset released by Meta AI and requires registration with Meta AI before download; it is not redistributed by this paper. Dataset details and the registration process are available at \url{https://ai.meta.com/datasets/segment-anything-video/}. All code implementing the baseline prompting policies, the P5 drift-detection and correction mechanism, and all twelve mitigation variants (A1, B1, A2, D1–D4, the three E1 pre-selection variants, the DAM4SAM-inspired multi-mask divergence check, and F1) is released alongside this paper at \url{https://github.com/bmnedashkivskyi-ai/no-oracle-drift-correction} (archived snapshot: Zenodo, DOI \url{https://doi.org/10.5281/zenodo.22811949}), together with the per-object result CSVs underlying every table in Section 4. All quantitative results in this paper were produced by the official SA-V evaluator (\texttt{sam2/sav\_dataset/sav\_evaluator.py}) and the official DAVIS evaluator (\texttt{davis2017-evaluation}), not by any custom or LLM-based scoring procedure.

\section*{Ethics Statement}

This work involves no human subjects research. It uses a publicly available (registration-gated) research benchmark, SA-V, for its intended purpose of video object segmentation evaluation, and a publicly released, frozen model checkpoint (SAM2.1). No personally identifying analysis was performed; SA-V videos may depict people, but this paper's contribution is a methodological finding about a drift-correction mechanism's aggregate and per-object segmentation accuracy, not any analysis of the individuals appearing in the footage.

\section*{Preregistration Statement}

This study was not preregistered. Standard guidance for preregistration (e.g., the exploratory-versus-confirmatory distinction underlying OSF-style registration practice) does not require it here: the work reports no confirmatory statistical hypothesis tests (Section 3.5), and the experimental series was iterative and exploratory by construction — each of the twelve mitigation attempts (Sections 3.4–4.5) was motivated by the measured outcome of the attempt before it, rather than specified in advance of data collection, and the further exploratory extension noted in Section 7 continues that same pattern rather than testing a separately pre-specified hypothesis. Readers should weigh the paper's findings accordingly: the distinction between what was anticipated and what was discovered along the way is preserved in the paper's own chronological framing of the results (Sections 3.4–4.5), rather than through a separate registered protocol.

\section*{AI-Use Disclosure}

This paper's writing was AI-assisted. The manuscript text (including this section) was drafted by an AI agent, working from a set of pre-existing, human-reviewed experimental result documents, with the writing task scoped to accurately transcribe and organize those results rather than to design or run experiments. The underlying experiments — the baseline prompting policy search, the P5 drift-correction mechanism and its calibration, and all twelve mitigation attempts — were designed and executed as part of an AI-agent-assisted empirical research project, in which an AI agent proposed hypotheses, implemented evaluation scripts, and ran experiments against the official SA-V and DAVIS evaluators under human oversight and decision-making at key branch points (e.g., whether to open the held-out test split, whether to proceed with the F1 direction despite a negative pilot, whether to pursue the DAM4SAM-inspired multi-mask divergence check as a literature-motivated follow-up). We state this plainly: both the experimental process and this write-up involved substantial AI agent involvement, and readers should weigh that accordingly alongside the fact that every reported number traces to an official-evaluator run on real model predictions rather than to any AI-generated or estimated figure.

\section*{Conflict of Interest}

The authors declare no conflict of interest.

\section*{Funding}

This research received no specific funding. No funding is declared.

\bibliographystyle{plainnat}
\bibliography{references}

\end{document}
